% !TeX root = ../main.tex
\chapter{Conclusion and Future Development}
\label{chap:ket-luan-huong-phat-trien}

\section{Achieved Results}

This thesis has proposed and built the foundation for an LMS-integrated
Generative AI system that automatically creates Micro-Content and Quizzes
from academic materials. The system does not stop at calling an LLM to
generate content; it also designs a knowledge-processing pipeline that
includes document extraction, chunking, embedding, retrieval, Knowledge
Graph, and Human-in-the-Loop review.

The main results are:
\begin{itemize}
    \item Designed the overall architecture including frontend, gRPC
          backend, workers, queues, and storage systems PostgreSQL,
          MinIO, Qdrant, and Neo4j.
    \item Applied Hexagonal Architecture/DDD to separate use cases from
          infrastructure adapters.
    \item Built the document-processing pipeline: parse, chunk, embed,
          store vectors, and generate AI content.
    \item Designed a Curriculum-Aware Layer to extract Learning Outcomes
          and map chunks to LOs.
    \item Oriented the LTI 1.3/OIDC integration so the system can operate
          as an external tool for an LMS.
    \item Built an evaluation framework for functionality, integration,
          retrieval, generated-content quality, and performance, with the
          principle that no metric is reported before it has actually
          been measured.
\end{itemize}

\section{System Limitations}

Although the system already has a relatively complete technical
foundation, several limitations still need to be addressed before
real-world deployment at classroom scale:
\begin{itemize}
    \item There has not yet been a large enough empirical evaluation
          across multiple courses and multiple instructors.
    \item Quiz quality depends on the LLM, prompts, and the quality of
          input chunks.
    \item OCR for scanned documents and documents containing many
          formulas has not yet reached production quality.
    \item GraphRAG is currently at the design/MVP curriculum-aware
          retrieval level and does not yet have a complete benchmark
          against pure RAG.
    \item Learning analytics dashboards, recommendation, and adaptive
          learning have not yet been completed as full product modules.
    \item Production deployment still requires additional monitoring,
          backup, secret management, and autoscaling.
\end{itemize}

\section{15-Week Phase-2 Implementation Plan (Graduation Thesis)}
\label{sec:plan-gd2}

Phase 2 of the project lasts 15 weeks with 4 members. Following the
supervisor's direction, the team is \textbf{not required to complete the
entire list} below. The minimum goal is to stabilize the end-to-end flow
and run a pilot with at least 10 students in one course, together with
empirical evaluation data for retrieval and quiz quality. Extension
items (advanced GraphRAG benchmark, advanced dashboard, multimodal
ingestion) are classified as \textbf{NICE-TO-HAVE}; if time is limited,
they may remain at design or prototype level without affecting the
conditions for defense.

\subsection{Four-Member Role Assignment}
\label{subsec:roles-gd2}

Roles are assigned according to the main \textit{runtime} each member is
responsible for, directly mapped to the 6 physical runtimes in
Table~\ref{tab:runtimes}. Each member acts as DRI (Directly Responsible
Individual) for their own scope while still cross-reviewing code from
other members to share knowledge.

\begin{longtable}{|p{1.2cm}|p{3.4cm}|p{9.5cm}|}
    \caption{Role and responsibility assignment for four members.}
    \label{tab:roles-gd2} \\
    \hline
    \textbf{ID} & \textbf{Main role} & \textbf{Responsibility scope} \\
    \hline
    \endfirsthead
    \hline
    \textbf{ID} & \textbf{Main role} & \textbf{Responsibility scope} \\
    \hline
    \endhead

    M1 & Backend AI (Python)
       & \texttt{document-worker}, \texttt{enrichment-worker},
         \texttt{chat-service}. Responsible for the parse $\to$ chunk
         $\to$ embed pipeline, Micro-Content/Quiz generation, RAG hybrid
         retrieval, and SSE streaming. Owner of Hexagonal ports
         \texttt{IParser}, \texttt{IEmbeddingModel},
         \texttt{ILLMClient}, \texttt{IVectorStore}. \\
    \hline
    M2 & Backend Core (NestJS) + Media (Go)
       & \texttt{core-api} (metadata, RBAC, Outbox poller),
         \texttt{video-worker} (FFmpeg, Whisper, segmentation), Neo4j
         sync, Curriculum ingestion, LTI AGS grade passback. Owner of
         the gRPC interface between runtimes. \\
    \hline
    M3 & Frontend / BFF (Next.js)
       & \texttt{web-bff}, instructor UI (review, dashboard,
         curriculum), learner UI (cards, quizzes, chatbot, video
         learning), LTI launch flow, mobile responsiveness. Owner of
         user experience and accessibility. \\
    \hline
    M4 & Data \& DevOps
       & Infrastructure (Docker Compose, secret management, backup),
         CI/CD, OpenTelemetry stack (Tempo, Loki, Grafana, Prometheus),
         schema migration, evaluation pipeline (Recall@K, MRR, quiz
         rubric), cost monitoring. Owner of runtime environments and
         evaluation. \\
    \hline
\end{longtable}

Two cross-cutting roles rotate weekly:
\begin{itemize}
    \item \textbf{Scrum Master:} coordinates daily standups, removes
          blockers, and updates the burndown chart. Rotates among
          M1--M4.
    \item \textbf{On-call pilot support:} supports instructors/students
          during the pilot period (Weeks 7--14). Rotates every 24h on a
          4-day cycle.
\end{itemize}

\subsection{Priority Classification and Minimum Success Criteria}
\label{subsec:priority-gd2}

Three priority levels are applied throughout the roadmap:
\begin{itemize}
    \item \textbf{[M] MUST:} mandatory for defense readiness -- stable
          end-to-end flow, minimum pilot, basic evaluation.
    \item \textbf{[S] SHOULD:} should be completed for a fuller report
          -- retrieval/quiz evaluation, basic instructor dashboard, cost
          monitoring.
    \item \textbf{[N] NICE-TO-HAVE:} implemented if time remains --
          advanced GraphRAG benchmark, advanced instructor dashboard,
          multimodal ingestion, Kubernetes manifests.
\end{itemize}

\textbf{Minimum success criteria (Definition of Done for the thesis):}
\begin{enumerate}
    \item LTI launch from Canvas into the system works for both
          \texttt{Instructor} and \texttt{Learner} roles.
    \item Upload PDF $\to$ parse $\to$ chunk $\to$ embed $\to$ Qdrant
          search returns relevant top-K results, with Recall@5 measured
          on a ground-truth set of $\geq 50$ queries.
    \item Micro-Content cards and Quiz items are generated from chunks,
          and instructors can review and publish them.
    \item The chatbot answers questions with citations to original
          chunks, and SSE streaming works inside the Canvas iframe.
    \item Video upload $\to$ STT $\to$ segment $\to$ Qdrant index, and
          students can watch video with timestamped transcripts.
    \item Pilot with $\geq 10$ students in at least 1 course, collecting
          $\geq 30$ quiz attempts and $\geq 50$ chat sessions.
    \item Report updates C5 (detailed implementation) and C6 (real
          pilot metrics, no remaining \texttt{TODO: measure empirically}).
\end{enumerate}

\subsection{15-Week Overview and 8 Sprints}
\label{subsec:overview-gd2}

\begin{longtable}{|p{1.4cm}|p{1.5cm}|p{4cm}|p{5.8cm}|}
    \caption{15-week Phase-2 overview divided into 8 sprints.}
    \label{tab:overview-gd2} \\
    \hline
    \textbf{Sprint} & \textbf{Weeks} & \textbf{Theme} & \textbf{Sprint
    starting milestone} \\
    \hline
    \endfirsthead
    \hline
    \textbf{Sprint} & \textbf{Weeks} & \textbf{Theme} & \textbf{Sprint
    starting milestone} \\
    \hline
    \endhead

    1 & 1--2 & Foundation, infra, LTI launch
      & Docker Compose can run the full stack; LTI launch demo works in
        Canvas sandbox. \\
    \hline
    2 & 3--4 & Document Pipeline + Micro-Content + Quiz
      & Upload PDF $\to$ chunk $\to$ Qdrant; draft card/quiz generation. \\
    \hline
    3 & 5--6 & Video Pipeline + RAG Chatbot + Video Search
      & Upload video $\to$ transcript; chatbot streams SSE with
        citations; Video Search returns YouTube + internal results. \\
    \hline
    4 & 7 & Hardening + Pilot kick-off
      & E2E stable for 48h; recruit pilot $\geq 10$ students; train the
        instructor. \\
    \hline
    5 & 8--9 & Pilot iteration + prompt tuning
      & Daily feedback-driven bug fixes; prompt iteration; performance
        profiling. \\
    \hline
    6 & 10--12 & Evaluation, security, polish
      & Recall@K/MRR have measured numbers; OWASP scan passes; UI
        polish. \\
    \hline
    7 & 13--14 & Pilot wrap-up + Data analysis
      & Export dataset; instructor/student survey; analyze quiz/chat
        quality. \\
    \hline
    8 & 15 & Final report + Demo
      & C5/C6 updated; demo video; defense rehearsal. \\
    \hline
\end{longtable}

\subsection{Overall Gantt Chart}
\label{subsec:gantt-gd2}

Figure~\ref{fig:gantt-sprints-gd2} shows the timeline of 8 sprints over
15 weeks with three main milestones (Pilot start -- week 7, Pilot end --
week 14, Submit + Defense -- week 15). Figure~\ref{fig:gantt-workstream-gd2}
details the plan by workstream and maps the main owner (M1--M4) for each
task, making each member's load and the cross-collaboration phases clear.

\begin{figure}[H]
\centering
\resizebox{\textwidth}{!}{%
\begin{ganttchart}[
    hgrid,
    vgrid={*{1}{draw=black!15}},
    x unit=0.95cm,
    y unit chart=0.55cm,
    y unit title=0.65cm,
    bar height=0.65,
    bar/.style={fill=blue!55, draw=blue!70},
    bar label font=\footnotesize,
    group/.style={fill=blue!80, draw=blue!90},
    group label font=\footnotesize\bfseries,
    milestone/.style={fill=red!80, draw=red!90},
    milestone label font=\footnotesize\itshape,
    title label font=\small\bfseries,
    title/.style={fill=gray!25, draw=gray!50},
    canvas/.style={draw=black!40}
]{1}{15}
\gantttitle{Phase-2 Week}{15} \\
\gantttitlelist{1,...,15}{1} \\
\ganttgroup{S1 Foundation + LTI}{1}{2} \\
\ganttgroup{S2 Document + Micro-Content}{3}{4} \\
\ganttgroup{S3 Video + RAG Chatbot}{5}{6} \\
\ganttgroup{S4 Hardening + Pilot kick-off}{7}{7} \\
\ganttgroup{S5 Pilot iter + Prompt tuning}{8}{9} \\
\ganttgroup{S6 Evaluation + Polish}{10}{12} \\
\ganttgroup{S7 Pilot wrap-up}{13}{14} \\
\ganttgroup{S8 Report + Demo}{15}{15} \\
\ganttmilestone{Pilot starts ($\geq$10 students)}{7} \\
\ganttmilestone{Pilot ends + freeze}{14} \\
\ganttmilestone{Submit + Defense}{15}
\end{ganttchart}%
}
\caption{Sprint-level Gantt chart for the 15-week Phase 2 with three
main milestones.}
\label{fig:gantt-sprints-gd2}
\end{figure}

\begin{figure}[H]
\centering
\resizebox{\textwidth}{!}{%
\begin{ganttchart}[
    hgrid,
    vgrid={*{1}{draw=black!15}},
    x unit=0.85cm,
    y unit chart=0.48cm,
    y unit title=0.6cm,
    bar height=0.65,
    bar/.style={fill=teal!55, draw=teal!70},
    bar label font=\scriptsize,
    group/.style={fill=teal!80, draw=teal!90},
    group label font=\scriptsize\bfseries,
    milestone/.style={fill=red!80, draw=red!90},
    milestone label font=\scriptsize\itshape,
    title label font=\small\bfseries,
    title/.style={fill=gray!25, draw=gray!50},
    canvas/.style={draw=black!40},
    progress label text={}
]{1}{15}
\gantttitle{Week}{15} \\
\gantttitlelist{1,...,15}{1} \\
% ===== M1: Backend AI (Python) =====
\ganttgroup{M1 Backend AI}{1}{15} \\
\ganttbar{Skeleton workers + Hexagonal Port}{1}{2} \\
\ganttbar{Parse + chunk + embed BGE-M3}{3}{4} \\
\ganttbar{Micro-Content + Quiz prompts}{3}{4} \\
\ganttbar{chat-service + RAG hybrid + RRF}{5}{6} \\
\ganttbar{Video Search MCP module}{5}{6} \\
\ganttbar{E2E smoke + bug fix}{7}{7} \\
\ganttbar{Prompt iteration from pilot}{8}{9} \\
\ganttbar{Retrieval eval (Recall@K, MRR)}{10}{12} \\
\ganttbar{Quiz/chat quality analysis}{13}{14} \\
\ganttbar{Update C5/C6 AI sections}{15}{15} \\
% ===== M2: Backend Core + Media =====
\ganttgroup{M2 Backend Core + Media}{1}{15} \\
\ganttbar{core-api NestJS + JWT verify}{1}{2} \\
\ganttbar{Outbox + Poller daemon}{3}{4} \\
\ganttbar{Neo4j sync + Curriculum extract}{3}{4} \\
\ganttbar{video-worker Go + FFmpeg + Whisper}{5}{6} \\
\ganttbar{LTI AGS + Rate limit}{7}{7} \\
\ganttbar{Curriculum versioning + LO migrate}{8}{9} \\
\ganttbar{Bug fix + hardening}{10}{12} \\
\ganttbar{Class diagram + C5 Core section}{13}{15} \\
% ===== M3: Frontend / BFF =====
\ganttgroup{M3 Frontend / BFF}{1}{15} \\
\ganttbar{web-bff + LTI login/launch + Canvas key}{1}{2} \\
\ganttbar{Upload UI + progress SSE + Review UI v1}{3}{4} \\
\ganttbar{Chat UI + Video learning + Card/Quiz}{5}{6} \\
\ganttbar{Onboarding flow + mobile QA}{7}{7} \\
\ganttbar{Instructor dashboard v1}{8}{9} \\
\ganttbar{UI polish + accessibility WCAG}{10}{12} \\
\ganttbar{Survey + final UI polish}{13}{14} \\
\ganttbar{Demo video + defense slides}{15}{15} \\
% ===== M4: Data \& DevOps =====
\ganttgroup{M4 Data + DevOps}{1}{15} \\
\ganttbar{Docker Compose + CI + Schema migrate}{1}{2} \\
\ganttbar{OpenTelemetry stack + integration test}{1}{4} \\
\ganttbar{Backup + Cost monitoring + Production env}{7}{7} \\
\ganttbar{Performance profiling}{8}{9} \\
\ganttbar{OWASP scan + Load test + Cost report}{10}{12} \\
\ganttbar{Export dataset + Cost aggregation}{13}{14} \\
\ganttbar{Appendix + RUNBOOK + release tag}{15}{15} \\
% ===== Cross-cutting =====
\ganttgroup{Cross-cutting}{1}{15} \\
\ganttbar{Daily standup + sprint review}{1}{15} \\
\ganttbar{Pilot operations (on-call rotation)}{7}{14} \\
\ganttbar{Incremental documentation}{2}{15} \\
\ganttmilestone{Pilot start}{7} \\
\ganttmilestone{Pilot stop}{14}
\end{ganttchart}%
}
\caption{Detailed Gantt chart by workstream and owner (M1--M4). Each
group corresponds to one member; the \emph{Cross-cutting} group contains
work shared by the whole team.}
\label{fig:gantt-workstream-gd2}
\end{figure}

Observations from the Gantt chart:
\begin{itemize}
    \item Weeks 1--2 have \textbf{uniformly high load} for all 4 members
          because scaffolding must be built in parallel. Pair programming
          in week 1 is recommended to share Hexagonal Port context
          between M1--M2.
    \item In weeks 5--6, M1 and M2 both carry heavy load (RAG +
          video-worker). M3 must finish base UI by week 4 so it does not
          create blockers.
    \item Week 7 is the \textbf{bottleneck}: hardening, pilot kick-off,
          and instructor/student training all happen together. It is the
          highest-risk week and should not include new feature work.
    \item In weeks 8--9, \textbf{M4 has lighter load} and can assist
          M1/M2 with curriculum versioning or run the evaluation pipeline
          early.
    \item Weeks 10--12 are the longest sprint, with \textbf{load
          concentrated on M1 and M4} for evaluation. M2 moves to
          hardening, while M3 polishes the UI.
    \item In weeks 13--15, the whole team shifts focus to
          \textbf{report + demo}; there are no new features.
\end{itemize}

\subsection{Sprint 1 (Weeks 1--2) --- Foundation and LTI Launch}
\label{subsec:sprint1-gd2}

Goal: build enough scaffolding so that all members can work in parallel
from Sprint 2. By the end of the sprint, the team must demonstrate LTI
launch from Canvas sandbox into \texttt{/manage/review} (instructor) or
\texttt{/learn/courses} (learner).

\begin{longtable}{|p{0.85cm}|p{5.35cm}|p{1.5cm}|p{0.9cm}|p{3.7cm}|}
    \caption{Sprint 1 --- Weeks 1--2.}
    \label{tab:sprint1-gd2} \\
    \hline
    \textbf{\#} & \textbf{Task} & \textbf{Owner} & \textbf{Prio} & \textbf{Acceptance} \\
    \hline
    \endfirsthead
    \hline
    \textbf{\#} & \textbf{Task} & \textbf{Owner} & \textbf{Prio} & \textbf{Acceptance} \\
    \hline
    \endhead

    1.1 & Setup monorepo + Docker Compose (15 containers) & M4 & M & \texttt{docker compose up} passes health checks. \\ \hline
    1.2 & Migrate PostgreSQL schema (24 + 3 tables from appendix) & M4 & M & \texttt{migrate up} runs cleanly on an empty DB. \\ \hline
    1.3 & CI workflow (GitHub Actions): lint + unit test & M4 & M & Opened PR passes CI under 10 minutes. \\ \hline
    1.4 & OpenTelemetry collector + Tempo + Loki + Grafana & M4 & S & Trace from \texttt{web-bff} to worker appears end-to-end in Grafana. \\ \hline
    1.5 & \texttt{core-api} skeleton (NestJS) + auth middleware (JWT verify) & M2 & M & \texttt{GET /health} and \texttt{GET /me} work after LTI launch. \\ \hline
    1.6 & \texttt{web-bff} skeleton + Next.js App Router & M3 & M & Homepage renders; routes \texttt{/lti/login}, \texttt{/lti/launch} reachable. \\ \hline
    1.7 & Canvas Developer Key + Cloud sandbox account & M3 & M & Tool can be registered in Canvas and launch test passes. \\ \hline
    1.8 & LTI 1.3 OIDC login + launch verification (JWKS) & M3 & M & Clicking an activity in Canvas $\to$ enters the system with correct role/context. \\ \hline
    1.9 & UPSERT \texttt{lms\_user\_mappings} + \texttt{lms\_course\_ref} after launch & M2 & M & Tables contain records after the first launch. \\ \hline
    1.10 & BFF-to-Core JWT signing (Ed25519 keypair) & M3, M2 & M & \texttt{core-api} rejects requests without a valid JWT. \\ \hline
    1.11 & \texttt{document-worker} skeleton + BGE-M3 model loading & M1 & M & Container starts and loads the model under 30s. \\ \hline
    1.12 & \texttt{enrichment-worker} skeleton + Gemini SDK adapter & M1 & M & Calls a sample prompt and parses strict JSON schema. \\ \hline
    1.13 & MinIO presigned URL upload flow (sample PDF) & M2, M3 & M & Browser uploads file $\to$ MinIO directly and passes. \\ \hline
    1.14 & Define Hexagonal Ports: \texttt{IParser}, \texttt{IEmbeddingModel}, \texttt{ILLMClient}, \texttt{IVectorStore}, \texttt{IGraphStore}, \texttt{IObjectStorage} & M1, M2 & M & ABC classes exist in codebase, with 2 adapters (real + mock). \\ \hline
\end{longtable}

\subsection{Sprint 2 (Weeks 3--4) --- Document Pipeline and Micro-Content}
\label{subsec:sprint2-gd2}

Goal: complete the document flow from upload to draft card/quiz
generation. By the end of the sprint, instructors can upload one sample
PDF, view its chunks in Qdrant, and review generated cards/quizzes
(publishing is not yet required).

\begin{longtable}{|p{0.85cm}|p{5.35cm}|p{1.5cm}|p{0.9cm}|p{3.7cm}|}
    \caption{Sprint 2 --- Weeks 3--4.}
    \label{tab:sprint2-gd2} \\
    \hline
    \textbf{\#} & \textbf{Task} & \textbf{Owner} & \textbf{Prio} & \textbf{Acceptance} \\
    \hline
    \endfirsthead
    \hline
    \textbf{\#} & \textbf{Task} & \textbf{Owner} & \textbf{Prio} & \textbf{Acceptance} \\
    \hline
    \endhead

    2.1 & PDF parser adapter (PyMuPDF) + OCR fallback (Tesseract/Docling) & M1 & M & Text-layer PDF parses; scanned PDF triggers OCR. \\ \hline
    2.2 & Heading-based chunker + size fallback (max 1500 tokens) & M1 & M & Chunks have correct \texttt{heading\_path} and \texttt{page\_range}. \\ \hline
    2.3 & BGE-M3 dense + sparse embedding + Qdrant upsert (Named Vectors) & M1 & M & Qdrant search returns top-K with both vectors. \\ \hline
    2.4 & Hybrid retrieval + RRF (preparation for Sprint 3) & M1 & S & RRF unit test merges 2 rankings correctly. \\ \hline
    2.5 & Document upload UI + progress SSE & M3 & M & Upload PDF $\to$ realtime progress bar visible. \\ \hline
    2.6 & Outbox table + Outbox Poller daemon in \texttt{core-api} & M2 & M & Insert outbox event $\to$ poll $\to$ enqueue BullMQ. \\ \hline
    2.7 & Neo4j sync worker: create \texttt{Chunk}, \texttt{SUPPORTS} & M2 & S & After ingesting 1 document, Neo4j has \texttt{Chunk} nodes. \\ \hline
    2.8 & Curriculum ingestion (syllabus PDF $\to$ Course, Chapter, LO, Assessment) using Gemini Flash JSON schema & M1, M2 & S & Ingesting a sample syllabus returns correct structure and inserts PostgreSQL rows. \\ \hline
    2.9 & \texttt{HeuristicLoMapper}: map chunk $\to$ LO by embedding similarity & M1 & S & One chunk is linked to $\geq 1$ \texttt{lo\_id} with confidence. \\ \hline
    2.10 & Micro-Content generation prompt + strict JSON schema & M1 & M & Card has \texttt{title}, \texttt{bullets}, \texttt{key\_insight}, \texttt{source\_chunk\_ids}. \\ \hline
    2.11 & Quiz generation prompt (Bloom level, mix of MCQ\_SINGLE / MCQ\_MULTI / TRUE\_FALSE / FILL\_BLANK) + validator & M1 & M & Quiz passes schema and does not contain ``All of the above''. \\ \hline
    2.12 & Review UI v1: list draft, view detail, approve/reject & M3 & M & Instructor can review 10 draft cards/quizzes. \\ \hline
    2.13 & State machine \texttt{GENERATED\_DRAFT} $\to$ \texttt{REVIEWING} $\to$ \texttt{APPROVED} & M2, M3 & M & Transitions are correct and audit log records them. \\ \hline
    2.14 & Integration test setup (testcontainers Postgres + Qdrant + Redis) & M4 & S & Document-pipeline test runs in CI under 5 minutes. \\ \hline
\end{longtable}

\subsection{Sprint 3 (Weeks 5--6) --- Video Pipeline and RAG Chatbot}
\label{subsec:sprint3-gd2}

Goal: complete the remaining two major AI subsystems -- video processing
and chatbot. By the end of the sprint, students can upload one lecture
video and ask questions through the AI Tutor chatbot with citations to
both documents and video segments.

\begin{longtable}{|p{0.85cm}|p{5.35cm}|p{1.5cm}|p{0.9cm}|p{3.7cm}|}
    \caption{Sprint 3 --- Weeks 5--6.}
    \label{tab:sprint3-gd2} \\
    \hline
    \textbf{\#} & \textbf{Task} & \textbf{Owner} & \textbf{Prio} & \textbf{Acceptance} \\
    \hline
    \endfirsthead
    \hline
    \textbf{\#} & \textbf{Task} & \textbf{Owner} & \textbf{Prio} & \textbf{Acceptance} \\
    \hline
    \endhead

    3.1 & \texttt{video-worker} (Go) skeleton + BullMQ Go client & M2 & M & Container can consume \texttt{VideoUploaded} events. \\ \hline
    3.2 & FFmpeg audio extraction subprocess + progress reporting & M2 & M & Extract 16kHz WAV from MP4/MOV. \\ \hline
    3.3 & Whisper STT integration (local CUDA + API fallback adapter) & M2 & M & Transcript has correct timestamps for a 5-minute video. \\ \hline
    3.4 & Silence + topic-shift segmentation, cut 30--90s clips & M2 & S & A 60-minute video creates $\geq 30$ segments. \\ \hline
    3.5 & Insert \texttt{video\_segments}, \texttt{transcript\_segments} + publish \texttt{TranscriptReady} & M2 & M & Event is received by \texttt{document-worker}. \\ \hline
    3.6 & Embed transcript segments into the same Qdrant collection \texttt{learning\_chunks} & M1 & M & Qdrant search returns both document chunks and video segments. \\ \hline
    3.7 & \texttt{chat-service} (FastAPI) skeleton + SSE endpoint & M1 & M & \texttt{/chat/stream} receives a query and returns token stream. \\ \hline
    3.8 & RAG hybrid retrieval (dense+sparse + RRF) + hard \texttt{course\_id} filter & M1 & M & Question for course A does not return chunks from course B. \\ \hline
    3.9 & Graph-aware reranking (boost chunks that \texttt{SUPPORTS} target LO) & M1 & S & Ablation exists: graph off $\to$ graph on, difference recorded. \\ \hline
    3.10 & Single-Use Ticket for SSE handshake & M2, M1 & M & Token is valid once and expires after 60s. \\ \hline
    3.11 & Citation extraction + response format \texttt{[1]}, \texttt{[2]} & M1 & M & Clicking citation opens the original chunk. \\ \hline
    3.12 & Store \texttt{chat\_messages} (UUIDv7) with citation metadata & M1 & M & Chat history reloads in correct order. \\ \hline
    3.13 & Video Search module: YouTube MCP + internal Qdrant transcript & M1 & S & Query ``Explain B-tree by video'' $\to$ returns 3 videos. \\ \hline
    3.14 & Chat UI: bubbles + citations + streaming caret & M3 & M & Token stream is smooth, without lag. \\ \hline
    3.15 & Video learning UI: player + transcript click-to-jump + related card panel & M3 & S & Clicking a transcript timestamp jumps the player correctly. \\ \hline
    3.16 & Card view + Quiz view (mobile-first) & M3 & M & Swipe gesture works; touch target $\geq 44$px. \\ \hline
\end{longtable}

\subsection{Sprint 4 (Week 7) --- Hardening and Pilot Kick-off}
\label{subsec:sprint4-gd2}

Goal: stabilize the system for 48 continuous hours and start the pilot
with a real class. This is a short sprint (1 week) but high risk, so
features should be frozen and only bug fixes should be accepted.

\begin{longtable}{|p{0.85cm}|p{5.35cm}|p{1.5cm}|p{0.9cm}|p{3.7cm}|}
    \caption{Sprint 4 --- Week 7.}
    \label{tab:sprint4-gd2} \\
    \hline
    \textbf{\#} & \textbf{Task} & \textbf{Owner} & \textbf{Prio} & \textbf{Acceptance} \\
    \hline
    \endfirsthead
    \hline
    \textbf{\#} & \textbf{Task} & \textbf{Owner} & \textbf{Prio} & \textbf{Acceptance} \\
    \hline
    \endhead

    4.1 & E2E smoke test: upload doc + video + generate card/quiz + chat + take quiz & All & M & One end-to-end flow passes within 30 minutes. \\ \hline
    4.2 & Backup scripts for PostgreSQL + Qdrant + Neo4j + MinIO & M4 & M & Restore drill: drop DB $\to$ restore $\to$ pass smoke test. \\ \hline
    4.3 & Cost monitoring dashboard (Prometheus rule + Grafana panel) & M4 & S & Token consumption is shown in real time by provider/use case. \\ \hline
    4.4 & Rate limit + quota enforcement (\texttt{user\_llm\_quota}) & M2 & M & User exceeding 100K tokens/day is rejected with 429. \\ \hline
    4.5 & LTI AGS grade passback (quiz score $\to$ Canvas Gradebook) & M2 & M & Student submits quiz $\to$ Canvas shows score. \\ \hline
    4.6 & Instructor onboarding flow (3 steps: connect Canvas, upload doc, review card) & M3 & M & New instructor completes setup under 15 minutes. \\ \hline
    4.7 & Mobile responsive QA pass (Android + iOS) & M3 & S & Card, quiz, chat usable on screens $\leq 375$px. \\ \hline
    4.8 & Recruit pilot: 1 course, 1 instructor, $\geq 10$ students & All & M & List and signed commitment exist. \\ \hline
    4.9 & Pilot user training session (1h, recorded video) & M3, M1 & M & Instructor + students know how to upload, review, learn, and chat. \\ \hline
    4.10 & Production env (VPS or small cloud) with HTTPS Caddy & M4 & M & Valid HTTPS domain, Canvas can launch. \\ \hline
    4.11 & Basic alert rules (error rate, queue depth, disk usage) & M4 & S & Discord/Email alert fires when rule triggers. \\ \hline
    4.12 & On-call rotation schedule (4 days/person) & All & M & Schedule shared with the supervisor. \\ \hline
\end{longtable}

\subsection{Sprint 5 (Weeks 8--9) --- Pilot Iteration and Prompt Tuning}
\label{subsec:sprint5-gd2}

Goal: keep the pilot stable for 2 weeks and collect enough data for the
evaluation in Sprint 6. The focus is feedback-driven bug fixing and
prompt refinement to improve card/quiz quality.

\begin{longtable}{|p{0.85cm}|p{5.35cm}|p{1.5cm}|p{0.9cm}|p{3.7cm}|}
    \caption{Sprint 5 --- Weeks 8--9.}
    \label{tab:sprint5-gd2} \\
    \hline
    \textbf{\#} & \textbf{Task} & \textbf{Owner} & \textbf{Prio} & \textbf{Acceptance} \\
    \hline
    \endfirsthead
    \hline
    \textbf{\#} & \textbf{Task} & \textbf{Owner} & \textbf{Prio} & \textbf{Acceptance} \\
    \hline
    \endhead

    5.1 & Daily standup + bug triage (15 minutes every morning) & All & M & 24h bug-free SLA for blocker errors. \\ \hline
    5.2 & Pilot bug fixing from feedback (rolling) & All & M & Pilot bug burndown chart decreases linearly. \\ \hline
    5.3 & Prompt iteration based on instructor quiz-quality feedback & M1 & S & V2 prompt reduces reject rate below 30\%. \\ \hline
    5.4 & Chat session quality audit (sample 20 sessions) & M1 & S & Hallucination rate is measured and mitigations exist. \\ \hline
    5.5 & Performance profiling: identify bottlenecks (Qdrant search, LLM call, parsing) & M4 & S & Report top-3 bottlenecks + tuning plan. \\ \hline
    5.6 & Curriculum versioning (track LO changes by \texttt{academic\_year}) & M2 & S & Creating LO v2 does not break old chunk links. \\ \hline
    5.7 & Instructor dashboard v1 (LO coverage, top wrong quiz items, class activity) & M3 & S & Basic dashboard exists, real-time not required. \\ \hline
    5.8 & Test on a second course (if a volunteer instructor is available) & All & N & One extended course experiment exists. \\ \hline
\end{longtable}

\subsection{Sprint 6 (Weeks 10--12) --- Evaluation, Security, and Polish}
\label{subsec:sprint6-gd2}

Goal: collect proper empirical metrics for Chapter 6, \textbf{replacing
all} \texttt{TODO: measure empirically} entries with real numbers. In
parallel, harden security and polish the UI.

\begin{longtable}{|p{0.85cm}|p{5.35cm}|p{1.5cm}|p{0.9cm}|p{3.7cm}|}
    \caption{Sprint 6 --- Weeks 10--12.}
    \label{tab:sprint6-gd2} \\
    \hline
    \textbf{\#} & \textbf{Task} & \textbf{Owner} & \textbf{Prio} & \textbf{Acceptance} \\
    \hline
    \endfirsthead
    \hline
    \textbf{\#} & \textbf{Task} & \textbf{Owner} & \textbf{Prio} & \textbf{Acceptance} \\
    \hline
    \endhead

    6.1 & Create retrieval ground truth: $\geq 50$ queries with correct chunks & M1 & M & Ground-truth JSONL file committed. \\ \hline
    6.2 & Measure Recall@1, Recall@5, Recall@10, MRR for 3 retrieval modes (dense, sparse, hybrid+RRF) & M1, M4 & M & Recall/MRR table in C6. \\ \hline
    6.3 & Measure nDCG for Graph-aware rerank (ablation on/off) & M1 & S & Comparison with/without Neo4j boost exists. \\ \hline
    6.4 & Quiz rubric: 5 criteria, instructor grades $\geq 100$ quiz items & M1, Supervisor & M & Rubric score distribution table in C6. \\ \hline
    6.5 & Chat quality: hallucination rate, citation hit rate, helpfulness Likert 5 & M1 & S & Result table in C6. \\ \hline
    6.6 & OWASP scan: \texttt{trivy}, \texttt{pip-audit}, \texttt{npm audit} & M4 & S & High vulnerabilities patched or whitelisted with reasons. \\ \hline
    6.7 & Prompt injection test set ($\geq 20$ scenarios) & M1 & S & Bypass rate below 10\%. \\ \hline
    6.8 & Load test (k6 or Locust) with 50 concurrent users & M4 & S & Chatbot p95 latency under 5s, error rate under 1\%. \\ \hline
    6.9 & UI polish: accessibility (WCAG AA), optional dark mode & M3 & S & Lighthouse accessibility $\geq 90$. \\ \hline
    6.10 & Cost report: total token usage + cost per pilot user & M4 & M & Cost table in C6. \\ \hline
    6.11 & Documentation: \texttt{README.md}, \texttt{ARCHITECTURE.md}, \texttt{RUNBOOK.md} & All & S & New developer onboarding under 1h. \\ \hline
    6.12 & GraphRAG ablation: compare with/without Neo4j boost (nDCG, MRR) & M1 & N & Comparison table in C6. \\ \hline
\end{longtable}

\subsection{Sprint 7 (Weeks 13--14) --- Pilot Wrap-up and Data Analysis}
\label{subsec:sprint7-gd2}

Goal: close the pilot, analyze the final data, and write results into
the report.

\begin{longtable}{|p{0.85cm}|p{5.35cm}|p{1.5cm}|p{0.9cm}|p{3.7cm}|}
    \caption{Sprint 7 --- Weeks 13--14.}
    \label{tab:sprint7-gd2} \\
    \hline
    \textbf{\#} & \textbf{Task} & \textbf{Owner} & \textbf{Prio} & \textbf{Acceptance} \\
    \hline
    \endfirsthead
    \hline
    \textbf{\#} & \textbf{Task} & \textbf{Owner} & \textbf{Prio} & \textbf{Acceptance} \\
    \hline
    \endhead

    7.1 & Stop pilot, export dataset (PostgreSQL dump + chat log JSONL + quiz attempt CSV) & M4 & M & Dataset has timestamps and is anonymized. \\ \hline
    7.2 & Instructor survey (10 Likert questions + 3 open questions) & M3 & M & $\geq 1$ instructor fully responds. \\ \hline
    7.3 & Student survey ($\geq 7$/10 students respond) & M3 & M & Quantitative and qualitative data collected. \\ \hline
    7.4 & Analyze quiz quality: manual scoring + LLM judge (Gemini Pro) comparison & M1 & S & Inter-annotator agreement Cohen's $\kappa$ computed. \\ \hline
    7.5 & Analyze chat sessions: top questions, top citations, failure cases & M1 & S & Pattern-analysis table in C6. \\ \hline
    7.6 & Aggregate cost (LLM tokens + Whisper + storage) & M4 & M & Cost per active user/month. \\ \hline
    7.7 & Final bug list + known issues for C6 Limitations section & All & M & Clear list in report. \\ \hline
    7.8 & Final UI polish (typos, copy, loading states) & M3 & M & Walkthrough finds no UI errors. \\ \hline
    7.9 & Freeze code backup, tag \texttt{v1.0-defense} & M4 & M & Tag exists on GitHub. \\ \hline
\end{longtable}

\subsection{Sprint 8 (Week 15) --- Final Report and Demo}
\label{subsec:sprint8-gd2}

Goal: finalize the report and prepare for defense. No new feature is
added.

\begin{longtable}{|p{0.85cm}|p{5.35cm}|p{1.5cm}|p{0.9cm}|p{3.7cm}|}
    \caption{Sprint 8 --- Week 15.}
    \label{tab:sprint8-gd2} \\
    \hline
    \textbf{\#} & \textbf{Task} & \textbf{Owner} & \textbf{Prio} & \textbf{Acceptance} \\
    \hline
    \endfirsthead
    \hline
    \textbf{\#} & \textbf{Task} & \textbf{Owner} & \textbf{Prio} & \textbf{Acceptance} \\
    \hline
    \endhead

    8.1 & Update C5 (Implementation): folder structure, adapter snippets, sample DDL/Cypher, Compose file & M1, M2, M3, M4 & M & C5 is $\geq 700$ lines and contains real code snippets. \\ \hline
    8.2 & Update C6 (Evaluation): replace \texttt{TODO} with pilot metrics, Recall/MRR/Rubric tables & M1, M4 & M & C6 has no \texttt{TODO}. \\ \hline
    8.3 & Update appendix: final schema, env config, prompt template & M4 & M & Appendix complete. \\ \hline
    8.4 & Draw real class diagram (\texttt{class\_diagram\_*.png}) & M2, M3 & S & C4 section 4.9 has figures. \\ \hline
    8.5 & Record 8--10 minute demo video (full instructor + learner journey) & M3 & M & Video uploaded to YouTube unlisted. \\ \hline
    8.6 & Defense slides (15--20 slides, 20-minute presentation) & All & M & Rehearsal passes before supervisor. \\ \hline
    8.7 & Internal defense rehearsal ($\geq 2$ times) & All & M & At least 10 prepared Q\&A questions. \\ \hline
    8.8 & Submit report by faculty deadline & All & M & Receipt available. \\ \hline
\end{longtable}

\subsection{Main Risks and Mitigations}
\label{subsec:risks-gd2}

\begin{longtable}{|p{4.5cm}|p{1.5cm}|p{8.5cm}|}
    \caption{Risks and mitigations.}
    \label{tab:risks-gd2} \\
    \hline
    \textbf{Risk} & \textbf{Level} & \textbf{Mitigation} \\
    \hline
    \endfirsthead
    \hline
    \textbf{Risk} & \textbf{Level} & \textbf{Mitigation} \\
    \hline
    \endhead

    Canvas Cloud sandbox is rate-limited or locked
        & High
        & Backup with self-hosted Canvas Open Source in Docker Compose;
          store credentials and sample course in advance. \\ \hline
    Lab GPU is unavailable for local Whisper
        & High
        & Use Whisper API adapter (OpenAI or AssemblyAI) as fallback;
          switch via env var \texttt{STT\_PROVIDER}. \\ \hline
    Gemini API rate limit or quota exhaustion
        & Medium
        & Configure Gemini Flash fallback; cache deterministic prompts;
          prepare $\geq 2$ backup keys. \\ \hline
    Pilot does not reach 10 students
        & High
        & Outreach from week 1; prepare $\geq 2$ backup instructors; two
          small classes can be combined. \\ \hline
    Pilot bugs spread and reduce student trust
        & High
        & Feature flags for quick rollback; 24h on-call during the first
          two pilot weeks; reset environment nightly if needed. \\ \hline
    LLM cost exceeds budget
        & Medium
        & Cost cap per user; Prometheus alerts; switch to Gemini Flash
          for tasks that do not require Pro. \\ \hline
    Privacy: internal course documents sent to external APIs
        & Medium
        & Hard filter PII before embedding; sign DPA with provider;
          provide a fully local option for sensitive courses. \\ \hline
    Member is sick or unexpectedly busy
        & Medium
        & Cross-review code from Sprint 1; pair programming in weeks 1--2
          to share context. \\ \hline
    Report is rejected because empirical data is missing
        & High
        & Sprint 6 is the longest sprint (3 weeks) dedicated to
          evaluation; replace all TODOs before Sprint 8. \\ \hline
\end{longtable}

\subsection{Plan Operation Notes}
\label{subsec:plan-notes-gd2}

\begin{itemize}
    \item \textbf{Daily standup} 15 minutes every morning, online via
          Discord or in person, recording blockers and assigning actions
          within 24h.
    \item \textbf{Sprint review + retro} at the end of each sprint with
          the supervisor, demoing achieved MUST tasks and adjusting the
          next sprint's priorities.
    \item \textbf{Mandatory code review:} every PR requires $\geq 1$
          reviewer who is not the branch owner. No self-merge except
          hot-fixes with a reason logged.
    \item \textbf{Flexible scope reduction:} if the pilot cannot start
          by the end of Sprint 4, move it to Sprint 5; if Recall@K is
          still unavailable by the end of Sprint 6, drop all
          NICE-TO-HAVE tasks to protect MUST criteria.
    \item \textbf{Report in parallel with code:} from Sprint 2 onward,
          every completed MUST feature must include 1--2 paragraphs of
          updates for C5 or C6 in the same PR, avoiding report-writing
          overload in Sprint 8.
\end{itemize}

In summary, Generative AI can significantly support the transformation
of academic materials into short learning content and assessment
questions, but it must be placed inside an architecture with source
traceability, instructor review, and clear evaluation. The approach that
combines RAG, Knowledge Graph, Bloom taxonomy, and LTI makes the system
more than a text-generation tool; it becomes a component that can be
integrated into the teaching and learning workflow of an LMS.
